\section{Equity: does it work for the patient in front of me?}
\label{sec:equity}

A clinician treats this patient, not the average patient. The fifth question asks
whether performance holds for the specific subgroup in the room, and whether bias
is actively surveilled rather than assumed away. The danger is well documented and
easy to miss, because a biased system can post excellent overall numbers while
failing a particular group. In a widely used health algorithm, reliance on
historical cost as a proxy for need systematically understated the illness of Black
patients, so that an aggregate that looked fair concealed a large disparity
\cite{obermeyer2019dissecting}. An average, in other words, is not evidence of
equity; it can be a place where inequity hides.

The framework answers this with subgroup performance reporting: the system's
accuracy, calibration, and error rates are measured and reported separately for
each pre-specified subgroup, not only in aggregate. This reporting is not a
one-time audit but a continuous obligation, because case mix and model behavior
drift over time. The ten-gate VVUQ examination is configured so that a subgroup in
which performance falls below threshold does not silently pass: the action for a
patient in an under-performing subgroup is held and routed to the ESCALATE path,
and the disparity itself becomes a surveillance priority rather than a footnote.

\autoref{tab:equity} sets out the subgroup axes the framework tracks, the metric
reported on each, and the action taken when a disparity appears. The axes span age,
ancestry, sex, comorbidity burden, and rural access, because a patient can be
disadvantaged on any of them and the relevant axis is rarely the one that is most
convenient to measure. Figure 7 shows the structure: overall performance is
decomposed by the subgroup audit, which resolves either to equitable performance or
to a flagged disparity that enters surveillance.

This is what lets equity be answered by measurement rather than by assurance. The
test a clinician can apply is direct: does it work for the patient in front of me,
specifically, on the axis that describes this patient? Because performance is
reported per subgroup and a shortfall triggers escalation and surveillance, the
clinician can ask that question and get a real answer instead of an aggregate that
flatters the system. Equity built this way also passes the family test, because the
patient in front of the clinician is treated as a member of a measured group, not
absorbed into an average that was never about them.

\begin{cfwfig}
\node[cfone] (a) at (0,0) {Overall performance}; \node[cfact] (b) at (4.2,0) {Subgroup audit};
\node[cfhope] (c) at (8.6,0.8) {Equitable}; \node[cfrisk] (d) at (8.6,-0.8) {Disparity flagged};
\draw[cfedge] (a)--(b); \draw[cfedge] (b)--(c); \draw[cfedge] (b)--(d);
\end{cfwfig}
\vspace{-0.2cm}
\cfwcaption{Figure 7. Subgroup equity (Mermaid 10 reproduced as colored TikZ):
overall performance is decomposed by a subgroup audit that resolves to equitable
performance or to a flagged disparity entering surveillance.}

\needspace{9\baselineskip}\begin{table}[H]
\footnotesize
\begin{tabularx}{\textwidth}{@{}L{3.2cm} Y L{3.2cm}@{}}
\toprule
\textbf{Subgroup axis} & \textbf{Metric reported} & \textbf{Action if disparity}\\
\midrule
Age & Accuracy and calibration by age band, including older patients & Escalate and add the band to surveillance\\
Ancestry & Error rates by reported ancestry to detect proxy bias & Escalate and audit the proxy variables\\
Sex & Sensitivity and specificity reported by sex & Escalate and recheck the training balance\\
Comorbidity burden & Performance stratified by comorbidity load & Escalate; widen the bound for complex cases\\
Rural access & Outcome and access gaps for remote sites & Escalate and route to the 24/7 response team\\
\bottomrule
\end{tabularx}
\caption{Subgroup audit axes: the metric reported on each axis and the action taken
when a disparity is detected.}
\label{tab:equity}
\end{table}

Equity is answered by measurement: performance reported per subgroup, with
disparities surfaced as surveillance priorities rather than buried in an average.
Writing per-subgroup reporting and bias surveillance into the illustrative
H. R. 9510 as a standing requirement \cite{Kawchak2026HR9510Bill} is the concrete
provision that lets a clinician trust the system works for the patient in front of
them, and not merely for the patient on average.
